% Part: set-theory
% Chapter: ordinals
% Section: introduction

\documentclass[../../../include/open-logic-section]{subfiles}

\begin{document}

\olfileid[hi]{sth}{ordinals}{intro}

\olsection{परिचय}

\olref[z][]{chap} में हमने \stagesinf{} के रूप में माना था कि पदानुक्रम का
एक अनंतवाँ चरण है (अनंतता का हमारा अभिगृहीत भी देखें)। किंतु
\stagessucc{} को मानकर हम अनंतवें चरण पर रुक नहीं सकते; हमें आगे बढ़ते रहना
होगा। अतः पहले अनंत चरण के बाद वाले अगले चरण पर हम उन समुच्चयों के सभी संभव
संग्रह बनाते हैं जो पहले अनंत चरण पर उपलब्ध थे; फिर दोहराते हैं; फिर
दोहराते हैं; फिर दोहराते हैं; \dots

यहाँ निहित रूप से यह हुआ है कि हमने संख्या की ऐसी ``सहज'' धारणा का आह्वान
आरंभ कर दिया है जिसके अनुसार समस्त प्राकृतिक संख्याओं के \emph{बाद} भी
संख्याएँ हो सकती हैं। विशेषतः यहाँ \emph{परिमितेतर क्रमसूचक} की धारणा
अंतर्भूत है। इस अध्याय का उद्देश्य इस विचार को अधिक कठोर बनाना है। हम
क्रमसूचक की सामान्य धारणा का अन्वेषण करेंगे और फिर कुछ समुच्चयों को स्पष्टतः
अपने क्रमसूचक परिभाषित करेंगे।

\end{document}
